% ============================================================
%  R. Nielsen, "Et Synspunkt for Darwinismen"
%  For Idé og Virkelighed, 1873, pp. 449–459
%  TRANSLATION (English)
% ============================================================
\documentclass[12pt,a4paper]{article}
\usepackage[utf8]{inputenc}
\usepackage[T1]{fontenc}
\usepackage[english]{babel}
\usepackage{microtype}
\usepackage{setspace}
\usepackage[margin=1.2in]{geometry}
\usepackage{fancyhdr}
\usepackage{hyperref}

\hypersetup{
  pdftitle={A Perspective on Darwinism},
  pdfauthor={R. Nielsen},
  hidelinks
}

\pagestyle{fancy}
\fancyhf{}
\rhead{Nielsen, \emph{A Perspective on Darwinism} (1873)}
\cfoot{\thepage}

\setstretch{1.3}

\title{\textbf{A Perspective on Darwinism}}
\author{R.\ Nielsen\\[4pt]
  \small Translated from the Danish}
\date{\emph{For Idé og Virkelighed}, 1873, pp.\ 449--459}

\begin{document}
\maketitle
\thispagestyle{fancy}

\bigskip

Darwin is not a natural philosopher, but---in the most distinguished
sense---a natural scientist. This simple observation, which does not
hint at any new discovery but merely states what anyone who concerns
himself with reviewing the Darwinian hypothesis can tell himself, is
in our judgment well suited to remove a portion of the
misunderstandings with which a fermenting public opinion seems, from
many quarters, to wish to encumber an attempt that is, in its
direction, magnificent.

The points around which the polemics of opinion revolve are:

\begin{enumerate}
  \item The common descent of organisms, and in particular the origin
        of the human being.
  \item The emergence of the first organic beings on earth---whether
        this is owed to a creation or merely to a distinctive
        co-operation of ordinary natural forces.
  \item The character of Darwinism as a whole---whether it can really
        be reconciled with human dignity, or must be said to contain
        tendencies that are hostile to the spirit and dangerous to
        religion and morality.
\end{enumerate}

For the clarification of these main points we have, as regards the
last in particular, absolutely no guidance in Darwin's purely personal
convictions, which are not more closely known to us and do not in any
case bear on the matter; we confine ourselves exclusively to the
theory as it stands before us.

\medskip

\noindent 1) ``Whence come these legions of living beings, these almost
numberless species and genera of plants and animals found in nature?''
Darwin answers: ``they all descend from a few simply organized forms;
they constitute one coherent system, whose individual members,
increasing in fullness and perfection, have in the course of time
worked their way forward in a sustained struggle for life.''

It will now easily be seen that it makes a great difference whether
we approach the view thus expressed from the standpoint of empirical
science or from that of natural philosophy. If one applies a
natural-philosophical measure, one is justified in asking for the
principle from which the theory as a whole can be derived.

In this respect not only philosophers, but natural scientists as well,
such as Th.\ L.\ Bischoff, have found themselves unsatisfied. ``The
Darwinian doctrine,'' says Bischoff, ``is no theory; it does not rest
on a principle which, because it is derived as a universally valid
truth, must also be acknowledged by all as a universally valid truth.
It contains many important and true thoughts and facts; but these can
only be acknowledged as correct insofar as they confirm themselves as
correct for objective experience and observation. But where this
verification fails or has not succeeded, we shall not be able to
attribute any significance to the Darwinian theory.''

Froschhammer,\footnote{\emph{Om Darwinismen}, translated by H.\
Steen. Copenhagen, 1873.} who judges Darwinism exclusively from a
philosophical-critical standpoint, lays weight, among other things,
on the obscurity that rests over Darwin's well-known catchword:
``natural selection'' (\emph{Qualit\"atsvalg}). Since some of his
other objections strike us as being of rather doubtful significance,
we shall single out this one in particular, which is, in a
natural-philosophical sense---mark this well---certainly substantially
grounded.

Froschhammer takes as his example the origin and development of the
eye. ``Since,'' he says, ``we must think of the primitive animal forms,
simple and imperfect as they are, as lacking eyes---just as there are
still animals in existence without eyes---the main question arises
afresh: how did eyes originally arise or how could they arise? They
must have arisen either by chance, or by an inexplicable,
incomprehensible \emph{generatio aequivoca}, or explicitly through a
new act of creation. In no case could they have arisen through natural
selection, since this, by its very concept, is capable only of
modifying, or properly speaking only of maintaining, what is already
given, and not of creating something new, something not yet given.''

The perfect eye is compared with the telescope, and the operation of
natural selection in perfecting the eye with the efforts of human
intelligence to improve that instrument; ``but,'' adds the
philosophical critic, ``surely wrongly. Unconscious nature is just as
little capable of imitating or performing the purposive activity of
opticians as it

is capable of imitating or replacing the activity of the artist (the
painter, or even merely the watchmaker). At this point Darwin, in
order not to come to a standstill in his explanation of the origin of
the most perfect eyes, resorts to a thoroughgoing personification of
natural selection. `Natural selection' is supposed to `watch closely'
and `carefully select,' and with never-failing tact to detect every
improvement that can promote further perfection.'' ``Were this to be
understood literally, Darwin himself would have introduced a
theological power into nature that would render all his other
explanatory attempts superfluous''---``and at the same time attributed
to natural selection a property that stands in flat contradiction with
its otherwise avowed character. But if the expressions named are to be
understood figuratively, then an explanation has been given or feigned
in words, where only something is figuratively asserted that cannot in
reality take place\ldots\ Natural selection cannot\ldots\ strive for
more perfect eyes, but can only maintain and make use of them once
they have in some way or other already come into being. And here,
then, there appears in a striking manner the very case of which Darwin
himself concedes that it could overthrow his theory.''

Similar objections can be raised against the manner in which the
origin of the human being (in the detailed work \emph{The Descent of
Man}) is explained. In order to prepare the reader for the coming
development, a human embryo is immediately in the beginning of the
first part compared with a dog embryo; from the similarity in outward
appearance an inference is drawn to similarity in nature---an
inferential procedure that is pervasive and characteristic of the
entire subsequent argument. Should anyone wish to throw himself with
critical zeal upon this point, it will scarcely be difficult for him
to discover and lay bare the fundamental weaknesses, once the matter
is, mark this well, treated philosophically.

But if we place ourselves at the standpoint of the researcher, the
whole affair takes on a different aspect. The natural scientist does
not have the task of constructing the typical peculiarities of living
beings and the differences actually obtaining among them from a
general principle; on the contrary! he finds the different forms given
to him in nature. His first question is not a question of necessity
but of possibility. If the many different species and genera were to
be assumed to stand independent of one another, then as many miracles
of creation would have to be presupposed as there are species; but to
build on miracles is to abandon all inquiry: whenever one is at a
loss to find a natural ground for a natural fact, one can in the most
convenient manner break off by saying: here a miracle has occurred.
``According to the ordinary view, that each individual being was the
result of its own creative act, we can only say that it is so, that
it has pleased the Creator to arrange all animals and plants in each
great class after a uniform plan; but this is not a scientific
explanation.''

These words of Darwin indicate clearly enough what it is he is aiming
at. Darwin's doctrine does not come forward as any finished
explanation and grounding of the peculiarity actually obtaining in
organic types, but as a hypothesis. This hypothesis has the great
scientific significance of leading inquiry into new tracks and
indicating a new method. Positing hypotheses is an easy matter; but
indicating new and fruitful methods is reserved for the

distinguished in science. What Darwin teaches us about the development
and transformation of organs in the struggle for life, about natural
selection, heredity, the laws of variation, and so forth, are decisive
determinations of the direction in which his investigations go---that
is, of his method. Instead of foisting upon him a natural-philosophical
principle that he does not have, one would do better to fix one's
entire attention on the distinctive character of the method, for the
researcher's principle is his method.

Darwin has not been refuted because he has not yet been able to show
how a mollusc has passed over to become either an arthropod or a
vertebrate, and still less how an ape has with philosophical necessity
become a human being; he aims, as far as we have understood his
method, not at the necessity of the forms that have come to be, but
at the possibility of gradual transitions and thereby at the
possibility of attaining ever greater insight into the natural inner
connection of organic forms with one another. As regards the
individual intermediate links, there is no one who has more
emphatically impressed upon us how immeasurably great our ignorance is
than Darwin himself.

\medskip

\noindent 2) How did the first organic beings come into existence on
earth? This question no natural scientist can answer, and it seems to
us that Darwin has upheld the standpoint of his science in the most
fitting manner by not giving any indication of wishing to answer it.
``There is,'' he says, ``grandeur in this view of life, that with its
several powers, having been originally breathed by the Creator into a
few forms or into one; and that, whilst this planet has gone cycling
on according to the fixed law of gravity, from so simple a beginning
endless forms most beautiful and most wonderful have been, and are
being, evolved.''

Truly a language worthy of a natural scientist! If the Darwinian
hypothesis is absolutely to be forced into a natural-philosophical
system and made to submit to a ``dialectical critique,'' it is
admittedly true that the researcher will find himself embarrassed. If
one has conceded one miracle, why not concede more? If the first
organic beings were created, why must not all of them have been
created? This is a critique that anyone can easily manage, and it is
therefore no wonder that it has been heard from so many quarters. But
what is the researcher's thought, properly speaking? He sees in life
an activity which, while it falls everywhere under universally valid
natural laws, is nonetheless an activity that cannot be derived or
explained from the interaction of elementary forces. The origin of
life is a riddle, and the riddle is intimated by saying that ``life
with its several powers was originally breathed by the Creator into a
few forms or into one.'' This reference to a creation fixes no
determinate miracle; it is left to each person to conceive of the
obscure beginning as he will: the origin and first beginning of life
is no object for empirical natural science.

With Darwin, the doctrine of creation stands in the same situation as
the doctrine of means and ends in nature (the teleological): creation
and purpose do not belong to the categories that find fruitful
application in empirical science. When the researcher finds himself
occasionally compelled to employ such designations, he is in a certain
embarrassment, because the words \emph{creation}, \emph{nature's
purpose}, and the like must be taken both in a figurative and a
literal sense at once: concerning actual creating as an

act, concerning a coming-into-being as ``a beginning of a being after
non-being,'' we have, says A.\ Humboldt, neither concept nor
experience. The same must be said of ``the purposive and yet
naturally necessary activity'' attributed by Darwin to natural
selection. That such fundamental concepts as creation and purpose
still bear a mystical obscurity is evidently philosophy's fault and
cannot be charged against empirical natural science. This leads us to
the consideration of the third main point.

\medskip

\noindent 3) Can Darwinism with justification be said to contain
tendencies that are hostile to the spirit and dangerous to religion
and morality? If it could be presupposed that the many who in the
present day have an imposing impression of the great, comprehensive
significance of natural science for cultural life also had a clear
understanding of what natural science is and must be as a science of
nature, it would be an easy matter to prove that Darwin's doctrine is
neither more nor less hostile to the spirit than any other
natural-historical theory.

Natural science is neither religious nor irreligious, neither ethical
nor unethical; if one wishes to attack Darwinism from a
religious-ethical starting-point, one simultaneously attacks the
entire far-branching natural science in all its members and parts.

What is supposed to be the offensive thing about Darwinism is that it
seeks to demonstrate the natural descent of higher species from lower,
because this is taken to be aimed at a demonstration of the natural
descent of the human being from some unknown animal form. It is, one
says, an affront to human dignity. Quite right!---when one reads into
it the religious idea of a supernatural act of creation and wants the
presuppositions of faith to be confirmed in science. Such a
confirmation one will admittedly not obtain from Darwin; but, one is
then to tell us, from which natural scientist one does expect to
obtain it.

Louis Agassiz, Darwin's declared adversary, describes his doctrine of
descent roundly as a scientific blunder, as untrue in its facts,
unscientific in its method, and pernicious in its tendency; Agassiz
holds strictly to the immutable limits of species and genera---can his
theory perhaps then be more religious than Darwin's? Agassiz proceeds
from zoogeography and draws from it consequences for the origin of the
human being. He speaks of creations and centres of creation and arrives
at no fewer than eight centres of creation with eight great zoological
realms: the Arctic, the Mongolian, the European, the American, the
Negro realm, the Hottentot realm, the Malayan, and the Australian
realm. Each of these realms has its characteristic fauna produced by a
local force, and its characteristic human race likewise produced by
the local force. Here then are creations enough. Far from coming into
ambiguous contact with animal transitional forms, the human being
stands independently over against the animal; for the human being was
created not merely once but eight times---that can count for something.
But that the eight-times-repeated creation, owed to the ``local
force'' present in eight places, has nothing to do with the religious
faith in the creation of the human being in God's image is obvious:
Agassiz's theory is just as little religious as Darwin's.

It is settled that natural science as a whole is neither religious nor
irreligious, neither pagan nor Christian, neither ethical nor
unethical; it is what it alone can and should be: objectively
inquiring science.

There are complaints, and not without grounds, that the cultivators of
science all too often misuse the results of inquiry to attack the
highest and most sacred truths of religion, and therewith of life; we
answer: when that is so, it is not science but the cultivators of
science that one must hold to account.

And what is it, then, that tempts the cultivators of science---whether
they be Darwinians or anti-Darwinians---to overreach? It is precisely
the old ingrained prejudice that religion's supernatural facts ought
to be capable of being grounded in science, that there really exists a
science built on miracles, a true miracle-science, which searches out
nature and history, the formation of the earth, the emergence of
living generations, the origin of the human being, and so forth, with
a quite different, deeper scientific gaze by means of miracles than
ordinary inquiry can manage without miracles. Whenever the alleged
miracle-science puts on its scientific air and sets out to correct the
researchers, its hollowness is exposed; and as it is sharply repulsed,
it not infrequently happens, unfortunately, that the higher truth
whose defence is conducted with false weapons is misjudged. So long
as one stubbornly insists on the old demand that science be religious
and religion scientific, the turning away from everything supernatural
and the denial of the word of revelation will continue to advance
steadily, for the mode of thinking of inquiry is as contrary as
possible to that of the Gospel.

Only when it is fully understood that religion has and must have its
own domain of knowledge, independent of science, and science its own
domains of investigation

independent of religion; only when, by the path of dearly bought
experiences, one has at last learned to see that it is just as
confounding to the spirit to wish to draw religious conclusions from
scientific premises as scientific conclusions from religious premises,
can there be any question of peace and conciliation between religion
and science. The only grounded objection that can be made against a
purported scientific theory is that it is unscientific. If Darwin's
method has scientific deficiencies, it will be corrected within
science, and natural scientists are the only proper, competent judges.
Whether it prevails or is defeated in the controversy has no essential
influence on the doctrine of the spirit.

The religious concept of creation is not a concept of knowledge but a
concept of faith. Darwinism does not harm religion, does not move
within faith's domain, does not resolve the problem of creation, does
not refute the revealed word: that the human being was created in
God's image.

\medskip
\begin{flushright}
R.\ Nielsen
\end{flushright}

\end{document}
